% Abstract — Part 1 focus (the instrument), with a one-line pointer to Part 2.
Evaluating large language models (LLMs) in regulated domains such as
biopharmaceutical compliance demands more than an aggregate accuracy
score: stakeholders need to know \emph{how hard} each question is and
\emph{where} a given system sits on a difficulty scale. We present an
automated pipeline that constructs a difficulty-calibrated question bank
for biopharma regulatory knowledge (FDA 21~CFR~Part~11, GCP, GMP,
promotional review, and informed consent) and calibrates it with
item-response theory (IRT). The pipeline (i)~generates compliance
scenarios under a ``Red Herring + subtle violation'' mandate that forces
genuine difficulty, (ii)~administers each item to a population of
synthetic AI examinees spanning models, sampling temperatures, and
system-prompt strictness, (iii)~grades responses with a dual-mode
LLM-as-judge, and (iv)~fits a Rasch (1PL) model, anchored to a fixed
baseline examinee, with classical item-fit quality control. Starting from
1{,}823 candidate items we freeze a healthy bank of
\textbf{1{,}284} items with calibrated difficulties spanning roughly
$\pm 4$ logits. A parallel track grounded in real FDA warning letters
confirms that enforced violations are systematically harder than
synthetic ones. In a second part, we treat the calibrated ability
$\theta$ as an \emph{optimization metric} and show that system-prompt
``strictness'' decomposes into a citation-\emph{format} effect and a
genuine \emph{reasoning} effect, which we separate with a strict/lenient
dual-judge design. We discuss the limits of a closed,
LLM-generated-and-judged evaluation loop and outline what external
validation would require. \todo{Tighten final numbers once py-irt
posterior SEs are wired in.}
